% ---------------------------------------------------------------------------
% LittleGPTracker - cheat sheet
%
% Build:
%   TEXMFVAR=$PWD/texmf/var TEXMFCONFIG=$PWD/texmf/config \
%       pdflatex -interaction=nonstopmode LittleGPTracker-Cheatsheet.tex
%
% TEXMFVAR *must* point at a writable directory.  The default
% (~/.texlive/texmf-var) is not writable here, and without it pdflatex cannot
% generate the EC bitmap fonts it needs and dies on the first \normalsize.
%
% pdflatex rather than lualatex/xelatex: this TeX Live has no Latin Modern
% OpenType fonts, so fontspec cannot load a default face.  Encoding stays at
% OT1 for the same reason - the T1/EC variants are bitmap only, while OT1 uses
% the Type1 CM fonts and stays sharp.  That means: no non-ASCII source, and
% <, > and | only inside \texttt.
% ---------------------------------------------------------------------------
\documentclass[8pt,a4paper]{extarticle}

\usepackage[a4paper,margin=8mm,top=7mm,bottom=7mm]{geometry}
\usepackage{multicol}
\usepackage{xcolor}
\usepackage{amssymb}
\usepackage{fancyhdr}
\usepackage[hidelinks]{hyperref}

% ------------------------------------------------------------------ palette
% The tracker's ayu dark theme, on paper: accents only, dark ink on white.
\definecolor{accent}{HTML}{A86400}   % ayu gold, darkened to stay legible
\definecolor{accent2}{HTML}{17604F}  % ayu teal, darkened
\definecolor{ink}{HTML}{10141C}
\definecolor{dim}{HTML}{545C6B}
\definecolor{rule}{HTML}{C9CDD6}
\definecolor{boxbg}{HTML}{F2EFE7}

\color{ink}
\pagestyle{fancy}
\fancyhf{}
\renewcommand{\headrulewidth}{0pt}
\fancyfoot[C]{\scriptsize\color{dim}LittleGPTracker cheat sheet \textbullet{} \thepage}

% ------------------------------------------------------------------- layout
\setlength{\parindent}{0pt}
\setlength{\parskip}{0.3mm}
\setlength{\columnsep}{4mm}
\setlength{\columnseprule}{0.2pt}
\renewcommand{\columnseprulecolor}{\color{rule}}

% Section heading: the text, then a rule that fills the rest of the line.  It
% used to be a full-width rule pulled up under the heading with a negative
% \vspace, which drew it straight through the descenders.
%
% \parfillskip is zeroed inside the group so the rule, not the paragraph's
% fill glue, takes all of the space left on the line.
\newcommand{\sechead}[1]{%
  \par\addvspace{1.5mm}%
  {\parfillskip=0pt\noindent{\small\bfseries\color{accent}#1}%
   \quad{\color{rule}\hrulefill}\par}%
  \addvspace{0.8mm}}

% Entries use an inline label with a small hanging indent, so a wrapped line
% lines up under the text rather than under the label.
\newcommand{\entry}[2]{\par\addvspace{0.28mm}\noindent\hangindent=2.4mm\hangafter=1 #1~#2\par}
\newcommand{\ctl}[2]{\entry{\texttt{\textbf{#1}}}{#2}}
\newcommand{\cmd}[3]{\entry{\textbf{\color{accent}#1}~{\scriptsize\texttt{#2}}}{#3}}
\newcommand{\sq}[1]{\texttt{#1}}
\newcommand{\kw}[1]{\textbf{#1}}

% shaded aside, kept to the column width
\newcommand{\note}[1]{%
  \par\addvspace{1.2mm}%
  \noindent\colorbox{boxbg}{%
    \begin{minipage}{\dimexpr\linewidth-2\fboxsep\relax}%
      \scriptsize\color{ink}#1%
    \end{minipage}}%
  \par\addvspace{0.6mm}}

\begin{document}

% ============================================================== masthead
\begin{center}
{\large\bfseries\color{accent} LittleGPTracker}\\[0.2mm]
{\scriptsize\color{dim}cheat sheet \textbullet{} 1.6.0-bacon20}
\end{center}
\vspace{-1.4mm}
{\color{rule}\hrule height 0.6pt}
\vspace{0.4mm}

\begin{multicols}{3}

\sechead{Buttons}
\kw{A}, \kw{B}, the shoulders \kw{LT} and \kw{RT}, \kw{Start} and
\kw{Select}. On a desktop build the same controls are Q and W for LT and RT,
the arrow keys for the D-pad and space for Start --- and the tricky part: the
\kw{A} key is \kw{B}, while \kw{S} is \kw{A}.

\ctl{D-pad}{Move the cursor.}
\ctl{A}{Insert note / chain / phrase.}
\ctl{A A}{Insert the next unused one.}
\ctl{A + D-pad}{Change the value under the cursor: up/down by \sq{10},
left/right by \sq{1}.}
\ctl{B}{Copy.}
\ctl{B + A}{Cut.}
\ctl{LT + A}{Paste.}
\ctl{B + D-pad}{Fast navigation: page, channel, or \sq{10}/\sq{1} in value
screens. On the song screen, left/right switches Song and Live mode.}
\ctl{RT + D-pad}{Walk between screens.}
\ctl{RT + B}{Mute or unmute the channel.}
\ctl{RT + A}{Solo the channel.}
\ctl{LT + RT}{Undo every mute and solo.}
\ctl{LT + up/down}{Jump to the next block after a blank row.}
\ctl{LT + (B,A)}{Clone the item into the next unused slot.}
\ctl{Start}{Play / stop, at the level of the screen you are on.}
\ctl{RT + Start}{Play the song from the edited row, on any screen.}

\sechead{Selections}
\ctl{LT + B}{Select what is under the cursor.}
\ctl{LT + B B}{Select the whole row.}
\ctl{LT + B B B}{Select the whole screen.}
\ctl{D-pad}{Grow or shrink it, then \kw{B} to copy or \kw{LT+A} to cut.}
\ctl{RT + B}{On a one-column selection in a chain or table: fill values upward.}

\sechead{Live mode}
\kw{B + left/right} on the song screen toggles it. \kw{Start} then queues only
the chain under the cursor, not the whole row, so each channel can be on a
different row.

\ctl{Start}{Queue; the play arrow blinks.}
\ctl{Start Start}{Queue immediately: starts at the end of the current phrase.}
\ctl{LT + Start}{Queue every channel on the row.}
\ctl{RT + Start}{Queue the channel to stop; the arrow becomes an underscore.}

\sechead{Screens}
\kw{RT + D-pad} walks the map below; the screen you are on is named top left.

{\scriptsize\ttfamily
\noindent project\\
song $\leftrightarrow$ chain $\leftrightarrow$ phrase $\leftrightarrow$ instrument\\
\hspace*{6mm}phrase $\updownarrow$ groove\\
\hspace*{6mm}phrase $\downarrow$ table \hspace*{3mm} instrument $\downarrow$ table\par}

\ctl{project}{Tempo, master, transpose, soft clip, housekeeping, \kw{midi} out,
render, save, and \kw{Exit} to quit the tracker.}
\ctl{song}{The order list: channels across, time down, a chain per cell.}
\ctl{chain}{Phrases, each row with its own transpose.}
\ctl{phrase}{Note, instrument, and two commands with their parameters.}
\ctl{instrument}{Sample, root note, loop, crush, filter, feedback, table.}
\ctl{table}{Three commands per row, run once per tick.}
\ctl{groove}{Ticks per step. Two sixes is straight; \sq{1F} swings.}

\sechead{Saving}
Stop playback first. Project screen, then \kw{Save Song}; the cursor vanishes
while it writes, so do not switch off until it returns. No undo, no autosave ---
save early and often.

\kw{Render} writes 16\,bit 44.1\,kHz WAVs as well as playing: \kw{Stereo} for the
mix, \kw{Stems} for a file per channel.

\sechead{Files}
Everything sits beside the binary.

\ctl{config.xml}{Screen, theme, MIDI input.}
\ctl{mapping.xml}{Buttons, joysticks and MIDI events, to tracker events.}
\ctl{lgpt\_*/}{A directory per project, holding \sq{lgptsav.dat}.}
\ctl{samplelib/}{The sample library.}
\ctl{log.txt}{Everything the tracker says, including the MIDI ports it found.}

\sechead{Commands}
Two per phrase row, three per table row. In \sq{aabb}, \sq{aa} is usually a
speed or tick count and \sq{bb} the target; speed \sq{00} is instant.

\cmd{ARPG}{abcd}{Arpeggio: cycles root, $+$\sq{a}, $+$\sq{b}, $+$\sq{c},
$+$\sq{d} semitones; a zero ends the cycle, so \sq{ARPG 3000} is root and $+3$.}
\cmd{CRSH}{aabb}{\sq{aa} pre-crush drive, \sq{00} leaves it alone; \sq{bb}
crush depth, \sq{0} is 1 bit, \sq{F} is 16 bit.}
\cmd{DLAY}{--bb}{Delay the note by \sq{bb} ticks.}
\cmd{FCUT}{aabb}{Cutoff to \sq{bb} at speed \sq{aa}: \sq{0080} jumps to half,
\sq{1000} closes it.}
\cmd{FLTR}{aabb}{Set cutoff \sq{aa} and resonance \sq{bb} outright.
\sq{FLTR 00FF} is unfiltered.}
\cmd{FRES}{aabb}{Resonance to \sq{bb} at speed \sq{aa}.}
\cmd{FBMX}{aabb}{Feedback mix to \sq{bb} at speed \sq{aa}.}
\cmd{FBTN}{aabb}{Feedback tune to \sq{bb} at speed \sq{aa}.}
\cmd{HOP}{aabb}{Jump to the next phrase in the chain at row \sq{bb};
\sq{bb}${=}$\sq{FF} stops the channel. Instant: the rest of the row still runs.}
\cmd{IRTG}{aabb}{Retrigger the instrument \sq{bb} semitones up. It accumulates:
\sq{IRTG 0001} climbs a semitone each time.}
\cmd{KILL}{--bb}{Stop the instrument after \sq{bb} ticks.}
\cmd{LEGA}{aabb}{Exponential slide to \sq{bb} at speed \sq{aa}; \sq{bb} is
relative, \sq{00}--\sq{7F} up and \sq{80}--\sq{FF} down.}
\cmd{LPOF}{aaaa}{Shift loop start and end by \sq{aaaa}, signed. Reset by the
next note.}
\cmd{MCHD}{aabb}{Two extra notes \sq{aa} and \sq{bb} semitones above the root,
held until the next note; \sq{00} is ignored.}
\cmd{MDCC}{aabb}{MIDI continuous controller \sq{aa} to value \sq{bb}, on the
running instrument's channel.}
\cmd{MDPG}{--bb}{MIDI program change; \sq{0000} is program 1.}
\cmd{MVEL}{--bb}{MIDI note velocity.}
\cmd{PAN}{aabb}{Pan to \sq{bb} at speed \sq{aa}.}
\cmd{PFIN}{aabb}{Pitch fine-tune to \sq{bb} at speed \sq{aa}; \sq{00} back to
the root.}
\cmd{PLOF}{aabb}{Sample offset. The sample is cut into 256 chunks: jump to
chunk \sq{aa}, or move \sq{bb} chunks relative.}
\cmd{PTCH}{aabb}{Linear pitch slide to \sq{bb} at speed \sq{aa}.}
\cmd{RTRG}{aabb}{Retrigger: loop \sq{bb} ticks from here, moving \sq{aa} ticks
forward on each pass.}
\cmd{TABL}{--bb}{Run table \sq{bb}.}
\cmd{TMPO}{--bb}{Set tempo. \sq{00} changes nothing, \sq{3C} is 60\,bpm,
\sq{190} is 400\,bpm.}
\cmd{VOLM}{aabb}{Volume to \sq{bb} at speed \sq{aa}. Start from an instrument
at volume 0 to fade sounds in.}

\sechead{Instrument}
\ctl{sample}{The wav. \kw{A A} opens the sample browser.}
\ctl{volume, pan}{Pan \sq{7F} is centre.}
\ctl{root note, detune}{Tuning. Your friend when an oscillator sounds wrong.}
\ctl{drive}{Gain before the crush.}
\ctl{crush}{Bit depth: \sq{F} clean, lower coarser.}
\ctl{downsample}{Halves the rate per step up; the aliasing whine.}
\ctl{flt cut/res}{Cutoff and resonance. Only a low-pass exists in this code.}
\ctl{type}{Mislabelled: it is wired to the filter \emph{mix}, not a filter type.}
\ctl{Mode}{Filter distortion: original, bassy, scream.}
\ctl{attenuate}{Volume after the filter.}
\ctl{fb tune, fb mix}{Feedback delay line. Below \sq{80} the length is the value,
above it ten times that.}
\ctl{interpolation}{\sq{linear} smooths; \sq{none} takes the nearest sample and
grits up low notes.}
\ctl{loop mode}{\sq{none}, \sq{loop}, \sq{ping pong}, \sq{oscillator} (the loop
becomes a tuned waveform), \sq{looper sync} (retunes the loop to exactly 16
rows).}
\ctl{slices}{Equal slices from C-2 up; \sq{1} disables.}
\ctl{start, loop start, loop end}{Hex sample offsets. End below start plays
backwards.}
\ctl{automation, table}{Automation steps the table one row per trigger instead
of letting it free-run.}

\kw{wet}/\kw{pad} print an impulse-response reverb onto the sample. It needs
ffmpeg; without it the row does nothing.

\sechead{MIDI}
Whether MIDI is available depends on the platform; the README carries a table.
Ports are listed once, at startup, in \sq{log.txt}:

{\scriptsize\ttfamily
\phantom{x}[RTMidiService] 1 input port(s)\\
\phantom{x}[RTMidiService] \phantom{1} Launchpad:0\par}

\kw{Input.} \sq{config.xml} holds \sq{MIDICTRLDEVICE}, matched as a
\emph{prefix} of the port name, so with no match no input is opened at all.
What an event then does is decided by \sq{mapping.xml}:

{\scriptsize\ttfamily
\phantom{x}\textless{}MAP src="midi:all:0:note:60"\\
\phantom{xxxxxx}dst="/event/start" /\textgreater\par}

That reads \sq{midi}:\emph{device}:\emph{channel}:\emph{type}:\emph{id}, where
the channel is \sq{0}--\sq{F} for MIDI 1--16, the device is a name from
\sq{log.txt} or \sq{all}, and the type is \sq{note}, \sq{cc}, \sq{pc}, \sq{pb}
or \sq{at}. Destinations are the ones the buttons use: \sq{/event/start},
\sq{/event/up}, \sq{/event/a} and so on.

\kw{Output.} Which interface the tracker sends to is a \emph{project} setting
--- the \sq{midi} line on the project screen --- and notes come from MIDI
instruments, whose screen sets channel, volume and gate length. While it plays
the tracker also sends clock, start and stop, so it can drive external gear.
The list is read when the project opens, so an interface plugged in later is
not offered until you reload.

\note{Ports are only enumerated at startup: connect the interface first, then
start the tracker. One plugged in later is not offered until you reload.}

\end{multicols}
\end{document}
